\section{Conclusions}
\label{sec:conclusions}

We have formulated an exact route from polarized amplitudes to analytic hard
coefficients for single-inclusive hadron production.  Its essential feature is
the preservation of physical information that ordinary scalar integral
manipulations can erase: positive-energy cuts, amplitude-side causal
prescriptions, BMHV evanescent tensors, and noninteger endpoint powers.

At NLO, six cut masters reduce to one elementary bubble and one generic top
function.  The top obeys a triangular first-order differential equation and
requires one non-elementary boundary period.  Its exact connection formula
determines the $x^{-\eps}$ and $x^{-1-2\eps}$ master modes.  After rational
hard coefficients are restored, the real contribution contains the two
physical sectors $(1-w)^{-1-\eps}$ and
$(1-w)^{-1-2\eps}$, from which the delta and plus distributions follow
without fitting endpoint logarithms.

For the NNLO double-real sector, exact IBP reduction yields 342 masters.  The
distinction between powered-integral equivalence, denominator equivalence,
differential closure, and equality of boundary periods reduces the
organizational problem without conflating it with the number of analytic
boundary constants.  The worked eight-master family demonstrates that a
non-Fuchsian coupled top sector and a nonuniform boundary can nevertheless be
solved analytically through $\eps^2$, with an independent numerical
comparison at every Laurent order.

The remaining task is concrete: construct the differential closures of the
17 maximal nontrivial NNLO denominator classes, determine their independent
physical boundary periods, evaluate all required masters analytically, and
assemble them with the exact rational coefficients.  Only then can the NNLO
endpoint distributions and the complete fixed-order hard functions be stated.
